%! Inverse Trigonometric Functions — Unit 3, Lesson 3.9 — Group Activity
\documentclass[10pt]{article}
\usepackage{precalculus-article}
\usepackage{precalculus-boxes}
\newcommand{\TallMath}[1]{$\displaystyle #1\rule[-1.4em]{0pt}{3.2em}$}

\begin{document}

\pageheader{Unit 3, Lesson 3.9}{Group Activity}
\namepartnerperiod

\vspace{0.10in}

% -----------------------------------------------------------------------
\begin{tcolorbox}[colback=white, colframe=black!40,
    title=\textbf{Tier R --- Remediate: Using a Table to Evaluate Inverse Trig},
    fonttitle=\bfseries, arc=2mm, left=3mm, right=3mm, top=2mm, bottom=2mm]

Use the reference table below to answer questions 1--6.

\begin{center}
\renewcommand{\arraystretch}{1.8}
\begin{tabular}{|c|c|c|c|}
\hline
$\theta$ & $\sin\theta$ & $\cos\theta$ & $\tan\theta$ \\
\hline
$-\dfrac{\pi}{2}$ & $-1$  & $0$                & undefined \\
$-\dfrac{\pi}{3}$ & $-\dfrac{\sqrt{3}}{2}$ & $\dfrac{1}{2}$ & $-\sqrt{3}$ \\
$-\dfrac{\pi}{4}$ & $-\dfrac{\sqrt{2}}{2}$ & $\dfrac{\sqrt{2}}{2}$ & $-1$ \\
$-\dfrac{\pi}{6}$ & $-\dfrac{1}{2}$ & $\dfrac{\sqrt{3}}{2}$ & $-\dfrac{\sqrt{3}}{3}$ \\
$0$ & $0$ & $1$ & $0$ \\
$\dfrac{\pi}{6}$ & $\dfrac{1}{2}$ & $\dfrac{\sqrt{3}}{2}$ & $\dfrac{\sqrt{3}}{3}$ \\
$\dfrac{\pi}{4}$ & $\dfrac{\sqrt{2}}{2}$ & $\dfrac{\sqrt{2}}{2}$ & $1$ \\
$\dfrac{\pi}{3}$ & $\dfrac{\sqrt{3}}{2}$ & $\dfrac{1}{2}$ & $\sqrt{3}$ \\
$\dfrac{\pi}{2}$ & $1$  & $0$  & undefined \\
$\dfrac{2\pi}{3}$ & $\dfrac{\sqrt{3}}{2}$ & $-\dfrac{1}{2}$ & $-\sqrt{3}$ \\
$\dfrac{3\pi}{4}$ & $\dfrac{\sqrt{2}}{2}$ & $-\dfrac{\sqrt{2}}{2}$ & $-1$ \\
$\pi$ & $0$ & $-1$ & $0$ \\
\hline
\end{tabular}
\end{center}

\textit{Remember:} $\arcsin$ outputs an angle in $[-\tfrac{\pi}{2},\tfrac{\pi}{2}]$;
$\arccos$ in $[0,\pi]$; $\arctan$ in $(-\tfrac{\pi}{2},\tfrac{\pi}{2})$.

\vspace{0.06in}
\begin{tabularx}{\linewidth}{@{} X X X @{}}
  \textbf{1.} $\arcsin\!\left(\dfrac{\sqrt{3}}{2}\right) = $ \blank{2cm}
  &
  \textbf{2.} $\arccos\!\left(\dfrac{1}{2}\right) = $ \blank{2cm}
  &
  \textbf{3.} $\arctan(\sqrt{3}) = $ \blank{2cm} \\[10pt]
  \textbf{4.} $\arcsin\!\left(-1\right) = $ \blank{2cm}
  &
  \textbf{5.} $\arccos(-1) = $ \blank{2cm}
  &
  \textbf{6.} $\arctan(0) = $ \blank{2cm}
\end{tabularx}
\end{tcolorbox}

\vspace{0.08in}

% -----------------------------------------------------------------------
\begin{tcolorbox}[colback=white, colframe=black!40,
    title=\textbf{Tier A --- Apply: Domain, Range, and Compositions},
    fonttitle=\bfseries, arc=2mm, left=3mm, right=3mm, top=2mm, bottom=2mm]

\textbf{7.} Complete the table.

\begin{center}
\renewcommand{\arraystretch}{2.0}
\begin{tabular}{|l|c|c|}
\hline
\textbf{Function} & \textbf{Domain} & \textbf{Range} \\
\hline
$f(x) = \arcsin x$    & \blank{3.5cm} & \blank{3.5cm} \\
\hline
$g(x) = \arccos x$    & \blank{3.5cm} & \blank{3.5cm} \\
\hline
$h(x) = \arctan x$    & \blank{3.5cm} & \blank{3.5cm} \\
\hline
\end{tabular}
\end{center}

\vspace{0.06in}
\textbf{8.} Evaluate $\sin\!\bigl(\arccos\tfrac{3}{5}\bigr)$.

\textit{Hint:} Let $\theta = \arccos\tfrac{3}{5}$.  Then $\cos\theta = \tfrac{3}{5}$.
Draw a right triangle and label the sides.

\writelines{3}

\vspace{0.04in}
\textbf{9.} Evaluate $\tan\!\bigl(\arcsin\tfrac{5}{13}\bigr)$.

\writelines{3}
\end{tcolorbox}

\vspace{0.08in}

% -----------------------------------------------------------------------
\begin{tcolorbox}[colback=white, colframe=black!40,
    title=\textbf{Tier E --- Extend: Compositions and Domain Analysis},
    fonttitle=\bfseries, arc=2mm, left=3mm, right=3mm, top=2mm, bottom=2mm]

\textbf{10.} Evaluate each expression.  Show all work.

\begin{tabularx}{\linewidth}{@{} X X @{}}
  \textbf{(a)} $\sin\!\bigl(\arcsin\tfrac{1}{2}\bigr) = $ \blank{2.5cm}
  &
  \textbf{(b)} $\arcsin\!\bigl(\sin\tfrac{\pi}{3}\bigr) = $ \blank{2.5cm} \\[8pt]
  \textbf{(c)} $\arcsin\!\bigl(\sin\tfrac{2\pi}{3}\bigr) = $ \blank{2.5cm}
  &
  \textbf{(d)} $\cos\!\bigl(\arccos(-0.7)\bigr) = $ \blank{2.5cm}
\end{tabularx}

\vspace{0.08in}
\textbf{11.} For which values of $x$ does $\sin(\arcsin x) = x$?  Explain.

\writelines{2}

\vspace{0.04in}
\textbf{12.} For which values of $x$ does $\arcsin(\sin x) = x$?  Explain
why the answer is not ``all real $x$.''

\writelines{3}

\vspace{0.04in}
\textbf{13.} Simplify $\cos(\arctan x)$ for $x > 0$.
Express your answer without any trig or inverse trig functions.
(\textit{Hint:} Let $\theta = \arctan x$; use a right-triangle diagram.)

\writelines{3}
\end{tcolorbox}

\end{document}
